\section{What must still become a four-dimensional quantum theory}
\label{pdh-existence-boundary}

The new coefficient removes a local ultraviolet obstruction. It does
not turn the analytic local recursion into the full Wilson measure.
The next construction must identify that recursion, with its retained
windows, gauge insertion, coupling extraction and vacuum convention,
inside an exact decomposition of the original finite Wilson integral.
Every mixed and large-field history must be included and estimated
with compatible boundary conditions. An arbitrary continuation of a
small-field germ outside its domain would define a different model.

The next priority is therefore a full-density matching theorem for
one fixed scheme. A useful statement would start from the actual
positive finite Wilson measure and express its pushforward as the
proved local sector plus explicitly controlled complementary sectors,
at the tuned couplings just constructed. It must hold uniformly as
the ultraviolet horizon grows. This statement is a target, not a
consequence of the one-loop calculation.

External observable sources then require their own uniform estimates.
Holomorphy in a connection background does not establish holomorphy
in a source coupled to a renormalized observable. Source counterterms,
normalization denominators, reflection-cell domains, and estimates
of all needed orders must be supplied for the same full measure.
Those estimates must produce a common continuum hierarchy with
compatible products, rather than separate limits of unrelated
observables.

Reflection positivity must pass from that same positive Wilson
sequence. Full Euclidean covariance, the required distributional
growth, a suitable vacuum sector, and reconstruction must be verified.
Finally, a nonzero scalar correlation alone does not identify
Yang--Mills: the limit needs an adequate mixed family of gauge-invariant
curvature observables and their short-distance normalization and
transformation laws. An interacting lower bound and the relevant
identification remain separate tasks. The mass-gap problem requires
an additional physical infrared argument.

The companion file
\texttt{work/final\_push\_existence\_dependency\_audit.txt}
records the detailed dependency map and source boundaries.
It explicitly distinguishes established working modules, conditional
reductions, and the open model-specific inputs. Earlier notebooks
remain chronological records; their older next-step paragraphs
must be read together with the results of the present notebook.

\begin{thebibliography}{9}
\bibitem{prior37}
\emph{Toward a four-dimensional Yang--Mills construction},
research Notebooks 15--36 in this workspace. In particular,
Notebook 32 supplies finite moment transfer and local shooting;
Notebook 33 fixes the continuum normalization and affine density
identification; Notebook 34 computes the exact unpinned scalar heat
coefficient; Notebooks 35--36 supply the actual mean-zero and soft
determinant identities and localized estimates.
\bibitem{vassilevich37}
D.~V.~Vassilevich, \emph{Heat kernel expansion: user's manual},
Physics Reports \textbf{388} (2003), 279--360,
\href{https://arxiv.org/abs/hep-th/0306138}{arXiv:hep-th/0306138}.
The reference coefficient is calculated explicitly here;
the source is not a proof of the actual lattice comparison.
\bibitem{adamslee37}
D.~H.~Adams and W.~Lee,
\emph{Structure of logarithmically divergent one-loop lattice
Feynman integrals}, Phys.\ Rev.\ D \textbf{77} (2008), 045010,
\href{https://arxiv.org/abs/0709.0781}{arXiv:0709.0781}.
Its fixed-integrand universality theorem is useful background;
uniformity in our varying effective actions is proved directly above.
\bibitem{jw37}
A.~Jaffe and E.~Witten, \emph{Quantum Yang--Mills Theory},
\href{https://www.claymath.org/wp-content/uploads/2022/06/yangmills.pdf}
{Clay Mathematics Institute problem description}.
\bibitem{source37}
E.~Sussman and collaborators,
\href{https://github.com/ethanwsussman/Yang--Mills-LLM-experiment}
{Yang--Mills-LLM-experiment}, deposited accounts pinned at commit
\texttt{e392ef45013a581664bab76b90867170b30fe10e}.
\end{thebibliography}
\end{document}
